\lyseis{1.4}{Ευκλείδεια διαίρεση, διαιρετότητα}
\sol{1} α) $1.752:24=73$ \quad β) $2.405:37=65$ \quad γ) $3.136:56=56$ \quad δ) $9.100:25=364$ \quad στ) $7.007:77=91$: όλες τέλειες, με υπόλοιπο 0 (δοκιμή, π.χ. $77\cdot 91=7.007$). ε) $8.514:42$ δίνει πηλίκο 202 και υπόλοιπο 30, γιατί $42\cdot 202=8.484$ και $8.514-8.484=30$. Δοκιμή: $42\cdot 202+30=8.514$.
\sol{2} α) $58=9\cdot 6+4$ \quad β) $200=7\cdot 28+4$ \quad γ) $365=30\cdot 12+5$ \quad δ) $1.000=64\cdot 15+40$ \quad ε) $47=50\cdot 0+47$ (το 50 δεν χωράει στο 47, πηλίκο 0) \quad στ) $0=12\cdot 0+0$.
\sol{3} α) Ναι, με διαιρέτη το 14 (πηλίκο 6): $\upsilon=6<14$. Όχι με διαιρέτη το 6, γιατί το υπόλοιπο 6 δεν είναι μικρότερο από το 6. β) Ναι, με διαιρέτη το 17 (πηλίκο 14) και με διαιρέτη το 14 (πηλίκο 17), γιατί $12<14$. γ) Όχι: το υπόλοιπο 25 δεν είναι μικρότερο από το 25, και είναι μεγαλύτερο από το 15. δ) Όχι: το υπόλοιπο 17 είναι μεγαλύτερο και από το 10 και από το 6.
\sol{4} Τα υπόλοιπα είναι μικρότερα από το 6: 0, 1, 2, 3, 4 ή 5.
\sol{5} $\Delta=11\cdot 48+7=528+7=535$. Δοκιμή: $535:11$ δίνει πηλίκο 48 και υπόλοιπο 7.
\sol{6} α) $\Delta=15\cdot 9+4=139$ \quad β) $250=12\cdot 20+10$, άρα $\pi=20$, $\upsilon=10$ \quad γ) $83=\delta\cdot 5+3$, άρα $5\delta=80$, $\delta=16$. Έλεγχος: $3<16$.
\sol{7} α) Λάθος: το υπόλοιπο είναι μικρότερο από τον διαιρέτη, άρα το πολύ 8. β) Σωστό. γ) Λάθος: δεν διαιρούμε με το 0. δ) Σωστό. ε) Λάθος: $23:1=23$. στ) Σωστό.
\sol{8} Η ισότητα ισχύει: $8\cdot 10+13=93$. Όμως το υπόλοιπο 13 δεν είναι μικρότερο από τον διαιρέτη 8: ο 8 χωράει άλλη μία φορά στο 13. Σωστό: $93=8\cdot 11+5$.
\sol{9} $214=52\cdot 4+6$: χρειάζονται 5 λεωφορεία. Στο πέμπτο υπάρχουν 6 παιδιά και μένουν $52-6=46$ κενές θέσεις. Με 4 λεωφορεία μένουν στο σχολείο 6 παιδιά.
\sol{10} $150=7\cdot 21+3$: Τρίτη και 3 μέρες, δηλαδή Παρασκευή. Πριν από 8 ημέρες: $8=7+1$, άρα μία μέρα πριν από την Τρίτη, δηλαδή Δευτέρα.
\sol{11} $500=24\cdot 20+20$: γεμίζει 20 αυλάκια και περισσεύουν 20 φυτά. Όχι, δεν φτάνουν για άλλο γεμάτο αυλάκι, γιατί $20<24$.
\sol{12} $87=13\cdot 6+9$: αγοράζει 6 μπάλες (78 €) και του μένουν 9 €.
\sol{13} $120:4=30$, $120:5=24$, $120:8=15$, $120:12=10$: τέλειες διαιρέσεις, γίνονται ίσες ομάδες. Το $120:7$ δίνει πηλίκο 17 και υπόλοιπο 1 ($120=7\cdot 17+1$) και το $120:9$ πηλίκο 13 και υπόλοιπο 3 ($120=9\cdot 13+3$): περισσεύουν 1 και 3 παιδιά αντίστοιχα.
\sol{14} $50=24\cdot 2+2$: μετά από 2 ολόκληρα εικοσιτετράωρα το ρολόι δείχνει πάλι 10:00 και μένουν 2 ώρες. Θα είναι 12:00, δύο μέρες αργότερα.
\sol{15} Οι αριθμοί είναι 32, 37, 42, 47, 52, 57 ($32=5\cdot 6+2$, $37=5\cdot 7+2$, \dots). Η διαφορά ανάμεσα σε δύο διαδοχικούς από αυτούς είναι πάντα 5, δηλαδή ο διαιρέτης.
\sol{16} Για $\nu=10$: $10:3$ έχει υπόλοιπο 1. Για 16: $16=3\cdot 5+1$, υπόλοιπο 1. Για 22: $22=3\cdot 7+1$, υπόλοιπο 1. Πάντα 1. Ο λόγος: $\nu=6\cdot\kappa+4$, το $6\kappa$ διαιρείται με το 3, και το 4 αφήνει υπόλοιπο 1.
\sol{17} Ο διαιρέτης πρέπει να είναι μεγαλύτερος από το υπόλοιπο 5, άρα ο μικρότερος είναι ο 6. Τότε $\Delta=6\cdot 9+5=54+5=59$.
\sol{18} Αν $6:0=x$, τότε $6=0\cdot x$, που είναι αδύνατο, γιατί το $0\cdot x$ είναι πάντα 0. Δεν υπάρχει $x$. Στο $0:0=x$ θα έπρεπε $0=0\cdot x$, που ισχύει για κάθε $x$: δεν υπάρχει ένας μοναδικός αριθμός. Και στις δύο περιπτώσεις η πράξη δεν ορίζεται.
\sol{19} $1.000.000=7\cdot 142.857+1$. Το υπόλοιπο είναι 1, άρα μία μέρα μετά την Τρίτη: Τετάρτη.
